参考资源
```Chantbook.tex
%========================================================
%  A5 Pali–English Chant booklet
%========================================================
\documentclass[12pt,a5paper]{article}

% 页面设置 ------------------------------------------------
\usepackage[
  top=15mm,
  bottom=15mm,
  left=8mm,
  right=6mm,
  headheight=10pt,
  headsep=6mm
]{geometry}


\usepackage{xcolor}
\usepackage{ragged2e}

% 标题设置 ------------------------------------------------
\usepackage{titlesec}
\definecolor{sectiongray}{gray}{0.35}

\titleformat{\section}
  {\large\bfseries\color{sectiongray}}
  {}
  {0pt}
  {}

\titlespacing*{\section}
  {0pt}
  {3ex}
  {1ex}

% XeLaTeX/LuaLaTeX 字体与语言 ------------------------------
% （polyglossia 会加载 fontspec，所以必须用 xelatex/lualatex）
\usepackage{polyglossia}
\setdefaultlanguage{english}
\setotherlanguage{pali}

\usepackage{fontspec}
% \setmainfont{Linux Libertine O}
\setmainfont{TeX Gyre Pagella}


% 微排版（XeLaTeX 下部分功能受限但可用） --------------------
\usepackage{microtype}

% 段落样式 ------------------------------------------------
\setlength{\parindent}{0pt}
\setlength{\parskip}{0.1\baselineskip}
\setlength{\emergencystretch}{1em}

% 页眉页码 ------------------------------------------------
\usepackage{fancyhdr}
\pagestyle{fancy}
\fancyhf{}
\fancyhead[C]{\normalsize\thepage}
\renewcommand{\headrulewidth}{0pt}
\renewcommand{\footrulewidth}{0pt}
\setlength{\headsep}{6mm}

% 链接（放在尽量后面） -----------------------------------
\usepackage[hidelinks]{hyperref}


% Pali / English 的样式 -----------------------------------
\newcommand{\palistyle}{\normalsize\bfseries\RaggedRight}
% \newcommand{\palistyle}{\normalsize\RaggedRight}
\newcommand{\englishstyle}{\small\normalfont\RaggedRight}

\newcommand{\chanttitle}[1]{%
  \par\bigskip
  \begin{center}
    {\fontsize{20}{24}\selectfont\bfseries #1}
  \end{center}
  \bigskip
}

%================ 自定义环境和命令 =======================
\newcommand{\chantgap}{%
  \par\vspace{0.6\baselineskip}%
}

\makeatletter

% 行计数器（每个小段内重置）
\newcounter{chantline}

% ===== 统一编号样式：把 1), 2) 的字体改小，并固定编号宽度 =====
\newlength{\chantlabelwidth}
\setlength{\chantlabelwidth}{1.0em} % 若超过 99) 可改 2.8em

\newcommand{\chantlabelnum}[1]{%
  \makebox[\chantlabelwidth][r]{\normalfont\footnotesize #1)}%
  % \makebox[\chantlabelwidth][l]{\normalfont\normalsize #1)}%
  \hspace{0.15em}%
}
\newcommand{\chantlabel}{\chantlabelnum{\thechantline}}
\newcommand{\chantlabelphantom}{%
  \makebox[\chantlabelwidth][l]{}%
  \hspace{0.15em}%
}

% 统一缩进（用于翻译/泰音译等对齐到正文起始）
\newcommand{\chantindent}{\chantlabelphantom}

% chantnote：不编号，但与 chantline 正文对齐
\newcommand{\chantnote}[1]{%
  {\palistyle\sloppy
    \chantindent #1\par
  }%
}

% chantnotetrans：与 chantline 正文对齐（用同一套缩进）
\newcommand{\chantnotetrans}[1]{%
  {\englishstyle\itshape\sloppy
    \chantindent #1\par
  }%
}

% 用来暂存本小段的 English 内容
\newcommand{\chant@english}{}

% 小段环境：一个 Pali+English 对应块
\newenvironment{chantsubsection}[1][]%
{%
  \par\smallskip
  \def\chant@title{#1}%
  \ifx\chant@title\@empty
  \else
    \noindent\textbf{\chant@title}\par\smallskip
  \fi
  \setcounter{chantline}{0}%
  \gdef\chant@english{}%
}%
{%
  \ifx\chant@english\@empty
  \else
    \par\vspace{0.4\baselineskip}%
    {\englishstyle\itshape
      \chant@english
    }\par\smallskip
  \fi
}

% -------- chantline：支持 “只写一行不带翻译” -----------
% 用法：
%   \chantline{Pali}{English}
%   \chantline{Pali}          % 不带翻译，不会在英文段落里出现空编号
\NewDocumentCommand{\chantline}{m g}{%
  \stepcounter{chantline}%
  {\palistyle\sloppy \chantlabel #1\par}%
  \IfNoValueF{#2}{%
    \begingroup
      \edef\tempb{\detokenize{#2}}%
      \ifx\tempb\@empty
        % { } 空翻译：不缓存
      \else
        \edef\temp{\endgroup
          \noexpand\g@addto@macro\noexpand\chant@english
            {\noexpand\chantlabelnum{\thechantline}\unexpanded{#2}\space}%
        }%
        \temp
      \fi
    \endgroup
  }%
}

% 不强制换行版：Pali 连成一段
\NewDocumentCommand{\chantlineinline}{m g}{%
  \stepcounter{chantline}%
  {\palistyle\sloppy \chantlabel #1\space}%
  \IfNoValueF{#2}{%
    \begingroup
      \edef\tempb{\detokenize{#2}}%
      \ifx\tempb\@empty
      \else
        \edef\temp{\endgroup
          \noexpand\g@addto@macro\noexpand\chant@english
            {\noexpand\chantlabelnum{\thechantline}\unexpanded{#2}\space}%
        }%
        \temp
      \fi
    \endgroup
  }%
}

% 双栏（每行两句 Pāli）+ 缓存一行 English
\newcommand{\chantlinepair}[3]{%
  \stepcounter{chantline}%
  {\palistyle\sloppy
    \chantlabel
    \begin{minipage}[t]{0.44\linewidth}#1\end{minipage}%
    \hspace{0.005\linewidth}%
    \begin{minipage}[t]{0.48\linewidth}#2\end{minipage}%
    \par
  }%
  \begingroup
    \edef\temp{\endgroup
      \noexpand\g@addto@macro\noexpand\chant@english
        {\noexpand\chantlabelnum{\thechantline}#3\space}%
    }%
  \temp
}

% 只输出原文（带编号），不缓存英文（保留给你备用）
\newcommand{\chantonly}[1]{%
  \stepcounter{chantline}%
  {\palistyle\sloppy \chantlabel #1\par}%
}
\newcommand{\chantonlyinline}[1]{%
  \stepcounter{chantline}%
  {\palistyle\sloppy \chantlabel #1\space}%
}

\makeatother

%================= Real footnotes style =================
\usepackage{footmisc}

% 1) 脚注上方的细横线：约 1/3 页面宽度
\renewcommand{\footnoterule}{%
  \kern -3pt
  \hrule width .33\linewidth height 0.4pt
  \kern 2.6pt
}

% 2) 单独控制脚注字号 + 不斜体（不影响你翻译的 \englishstyle）
\makeatletter
\renewcommand\footnotesize{%
  \@setfontsize\footnotesize{8.5pt}{10pt}% 这里改成你想要的脚注大小/行距
  \normalfont % 强制正体（非斜体）
}
\makeatother

%（可选）脚注段落格式更紧凑一点
\setlength{\footnotesep}{0.6\baselineskip}


%========================================================
\begin{document}
\RaggedRight

\chanttitle{Morning Chanting}

% 01-ratanattaya.tex
% Section [1] Ratanattaya-vandanā

\section*{m1 | Ratanattaya-vandanā}\label{sec:m01}

\begin{chantsubsection}
  \chantline
    {Yo so bhagavā arahaṃ sammāsambuddho}
    {That Blessed One is an arahant (a fully awakened one), perfectly awakened.}

  \chantline
    {Svākkhāto yena bhagavatā dhammo}
    {The Dhamma (the Buddha's teaching) is well expounded by that Blessed One.}

  \chantline
    {Supaṭipanno yassa bhagavato sāvakasaṅgho}
    {The Sangha (the community of the Blessed One's disciples) has practiced well.}

  \chantline
    {Tammayaṃ bhagavantaṃ sadhammaṃ sasaṅghaṃ}
    {To that Blessed One, to the true Dhamma, and to the Sangha}

  \chantline
    {Imehi sakkārehi yathārahaṃ āropitehi abhipūjayāma}
    {we offer these offerings, rightly arranged, and pay deep homage.}

  \chantline
    {Sādhu no bhante bhagavā suciraparinibbutopi}
    {It is well, Venerable Sir, that the Blessed One, though long since attained to final Nibbāna}

  \chantline
    {Pacchimājanatānukampamānasā}
    {out of compassion for later generations,}

  \chantline
    {Ime sakkāre duggatapaṇṇākārabhūte paṭiggaṇhātu}
    {please accept these humble, trifling offerings,}

  \chantline
    {Amhākaṃ dīgharattaṃ hitāya sukhāya}
    {for our long-lasting welfare and happiness.}
\end{chantsubsection}


%========================================
% Section：[2] Ratanattayanamakārapāṭha
%========================================
% title-en: Salutation to the Triple Gem
\section*{m2 | Ratanattayanamakārapāṭha}\label{sec:m02}

\begin{chantsubsection}
  \chantline
    {Arahaṃ sammāsambuddho bhagavā}
    {The Blessed One is an arahant, perfectly awakened;}

  \chantlineinline
    {Buddhaṃ bhagavantaṃ abhivādemi \; --- bow ---}
    {To the Buddha, the Blessed One, I pay homage.}

  \chantlineinline
    {Svākkhāto bhagavatā dhammo dhammaṃ namassāmi \; --- bow ---}
    {The Dhamma is well expounded by the Blessed One; to the Dhamma I pay homage.}

  \chantline
    {Supaṭipanno bhagavato sāvakasaṅgho saṅghaṃ namāmi}
    {The Sangha of the Blessed One’s disciples has practiced well; to the Sangha I pay homage.}
  \chantnote{--- bow ---}

\end{chantsubsection}

%========================================
% Section：[3] Pubbabhāganamakārapāṭha
%========================================
% title-en: Preliminary Salutation Verses
\section*{m3 | Pubbabhāganamakārapāṭha}\label{sec:m03}

\begin{chantsubsection}
  \chantnote
    {\{Handa mayaṃ buddhassa bhagavato
     pubbabhāganamakāraṃ karoma se\}}
  \chantnotetrans
    {\{Now let us pay preliminary homage to the Buddha, the Blessed One.\}}

  \chantnote
    {Namo tassa bhagavato arahato sammāsambuddhassa}
  \chantnote
    {--- three times ---}
  \chantgap
  \chantnotetrans
    {Homage to that Blessed One, an arahant, perfectly awakened.}
\end{chantsubsection}

%========================================
% Section：[4] Buddhābhithuti
%========================================
% title-en: Praise of the Buddha
\section*{m4 | Buddhābhithuti}\label{sec:m04}

\begin{chantsubsection}
  \chantnote
    {\{Handa mayaṃ buddhābhithutiṃ karoma se\}}
  \chantnotetrans
    {\{Now let us recite in praise of the Buddha.\}}
  \chantgap
  \chantline
    {Yo so tathāgato arahaṃ sammāsambuddho}
    {He is the Tathāgata (the Thus-Gone One), an arahant, perfectly awakened,}

  \chantline
    {Vijjācaraṇasampanno sugato lokavidū}
    {endowed with true knowledge and conduct, the Well-gone One, knower of the worlds,}

  \chantline
    {Anuttaro purisadammasārathi}
    {unsurpassed trainer of persons to be tamed,}

  \chantline
    {Satthā devamanussānaṃ buddho bhagavā}
    {teacher of gods and humans; the awakened, Blessed One.}

  \chantline
    {Yo imaṃ lokaṃ sadevakaṃ samārakaṃ sabrahmakaṃ}
    {He has realized this world with its gods, Māras, and Brahmās,}

  \chantline
    {Sassamaṇabrāhmaṇiṃ pajaṃ sadevamanussaṃ}
    {and this generation with its ascetics and brahmins, its gods and humans,}

  \chantline
    {Sayaṃ abhiññā sacchikatvā pavedesi}
    {and teaches it after having realized it by his own direct knowledge.}

  \chantline
    {Yo dhammaṃ desesi ādikalyāṇaṃ}
    {He teaches the Dhamma that is admirable in the beginning,}

  \chantline
    {Majjhekalyāṇaṃ pariyosānakalyāṇaṃ}
    {admirable in the middle, and admirable in the end,}

  \chantline
    {Sātthaṃ sabyañjaṇaṃ kevalaparipuṇṇaṃ parisuddhaṃ}
    {with the meaning and wording complete, perfect in every respect, wholly pure.}

  \chantline
    {Brahmacariyaṃ pakāsesi}
    {He has made known the perfectly complete holy life.}

  \chantline
    {Tamahaṃ bhagavantaṃ abhipūjayāmi}
    {To that Blessed One I pay deep homage.}

  \chantline
    {Tamahaṃ bhagavantaṃ sirasā namāmi \;\;\; --- bow ---}
    {To that Blessed One I bow my head.}

  % \chantnote{--- bow ---}
\end{chantsubsection}

%========================================
% Section：[5] Dhammābhithuti
%========================================
% title-en: Praise of the Dhamma
\section*{m5 | Dhammābhithuti}

\begin{chantsubsection}
  \chantline
    {Handa mayaṃ dhammābhithutiṃ karoma se}
    {Now let us recite in praise of the Dhamma.}

  \chantline
    {Yo so svākkhāto bhagavatā dhammo}
    {The Dhamma is well expounded by the Blessed One,}

  \chantline
    {Sandiṭṭhiko akāliko ehipassiko}
    {visible here and now, timeless, inviting one to come and see,}

  \chantline
    {Opanayiko paccattaṃ veditabbo viññūhi}
    {leading inwards, to be realized individually by the wise.}

  \chantline
    {Tamahaṃ dhammaṃ abhipūjayāmi}
    {To this Dhamma I pay deep homage.}

  \chantline
    {Tamahaṃ dhammaṃ sirasā namāmi}
    {To this Dhamma I bow my head.}
  \chantnote{--- bow ---}

\end{chantsubsection}

%========================================
% Section：[6] Saṅghābhithuti
%========================================
\section*{m6 | Saṅghābhithuti}

\begin{chantsubsection}
  \chantline
    {Handa mayaṃ saṅghābhithutiṃ karoma se}
    {Now let us recite in praise of the Sangha.}

  \chantline
    {Yo so supaṭipanno bhagavato sāvakasaṅgho}
    {The Sangha of the Blessed One’s disciples has practiced well,}

  \chantline
    {Ujupaṭipanno bhagavato sāvakasaṅgho}
    {has practiced uprightly,}

  \chantline
    {Ñāyapaṭipanno bhagavato sāvakasaṅgho}
    {has practiced in the right way,}

  \chantline
    {Sāmīcipaṭipanno bhagavato sāvakasaṅgho}
    {has practiced properly,}

  \chantline
    {Yadidaṃ cattāri purisayugāni aṭṭha purisapuggalā}
    {namely, the four pairs, the eight kinds of noble individuals.}

  \chantline
    {Esa bhagavato sāvakasaṅgho}
    {This is the Sangha of the Blessed One’s disciples:}

  \chantline
    {Āhuneyyo pāhuneyyo dakkhiṇeyyo añjalīkaraṇīyo}
    {worthy of gifts, worthy of hospitality, worthy of offerings, worthy of reverential salutation,}

  \chantline
    {Anuttaraṃ puññakkhettaṃ lokassa}
    {the unsurpassed field of merit for the world.}

  \chantline
    {Tamahaṃ saṅghaṃ abhipūjayāmi}
    {To this Sangha I pay deep homage.}

  \chantline
    {Tamahaṃ saṅghaṃ sirasā namāmi \;\;\; --- bow ---}
    {To this Sangha I bow my head.}

\end{chantsubsection}
%========================================
% Section：[7] Ratanattayappaṇāmagāthā
%========================================
\section*{m7 | Ratanattayappaṇāmagāthā}

\begin{chantsubsection}
    \chantnote{\{Handa mayaṃ ratanattayappaṇāmagāthāyo ceva saṃvegaparikittanapāṭhañca bhaṇāma se\}}
    \chantnotetrans{\{Now let us recite our verses of salutation to the Triple Gem and a passage to stir urgency.\}}
  \chantgap
  \chantline
    {Buddho susuddho karuṇāmahaṇṇavo}
    {The Buddha, utterly pure, with compassion deep as the ocean,}

  \chantline
    {Yoccantasuddhabbarañāṇalocano}
    {possessing the eye of excellent, stainless wisdom,}

  \chantline
    {Lokassa pāpūpakilesaghātako}
    {the destroyer of the world’s evil and defilements:}

  \chantline
    {Vandāmi buddhaṃ ahamādarena taṃ}
    {that Buddha I revere with a devoted heart.}

  \chantline
    {Dhammo padīpo viya tassa satthuno}
    {The Dhamma of that Teacher, like a lamp,}

  \chantline
    {Yo maggapākāmatabhedabhinnako}
    {distinguished as the Path, Fruition, the Deathless and its divisions,}

  \chantline
    {Lokuttaro yo ca tadatthadīpano}
    {supramundane, illuminating its own objective,}

  \chantline
    {Vandāmi dhammaṃ ahamādarena taṃ}
    {that Dhamma I revere with a devoted heart.}

  \chantline
    {Saṅgho sukhettābhayatikhettasaññito}
    {The Sangha, known as a fertile field and a field of fearlessness,}

  \chantline
    {Yo diṭṭhasanto sugatānubodhako}
    {the seers who realize, awakened in the footsteps of the Well-Gone One,}

  \chantline
    {Lolappahīno ariyo sumedhaso}
    {free from fickleness, noble and of excellent wisdom:}

  \chantline
    {Vandāmi saṅghaṃ ahamādarena taṃ}
    {that Sangha I revere with a devoted heart.}

  \chantline
    {Iccevamekantabhipūjaneyyakaṃ}
    {Thus this threefold basis, supremely worthy of veneration,}

  \chantline
    {Vatthuttayaṃ vandayatābhisaṅkhataṃ}
    {when honored and attended to in this way,}

  \chantline
    {Puññaṃ mayā yaṃ mama sabbupaddavā}
    {may the merit I have made cause all my misfortunes}

  \chantline
    {Mā hontu ve tassa pabhāvasiddhiyā}
    {not to come to be, through the fulfillment of its power.}
\end{chantsubsection}

%========================================
% Section：[8] Saṃvegaparikittanapāṭha
%========================================
\section*{m8 | Saṃvegaparikittanapāṭha}

% ---------- 第一段 ----------
\begin{chantsubsection}
  \chantlineinline
    {Idha tathāgato loke uppanno arahaṃ sammāsambuddho,}
    {Here in this world the Tathāgata, an arahant, perfectly awakened, has arisen.}
  \chantlineinline
    {Dhammo ca desito niyyāniko upasamiko parinibbāniko
     sambodhagāmī sugatappavedito.}
    {And the Dhamma has been taught that leads out, brings peace, leads to final Nibbāna,
     leads to awakening, proclaimed by the Well-gone One.}
  \chantlineinline
    {Mayantaṃ dhammaṃ sutvā evaṃ jānāma.}
    {Having heard that Dhamma, we understand thus:}
  \chantlineinline
    {Jātipi dukkhā, jarāpi dukkhā, maraṇampi dukkhaṃ,}
    {birth is suffering, aging is suffering, death is suffering;}
  \chantlineinline
    {sokaparidevadukkhadomanassupāyāsāpi dukkhā,}
    {sorrow, lamentation, pain, grief and despair are suffering;}
  \chantlineinline
    {Appiyehi sampayogo dukkho, piyehi vippayogo dukkho,}
    {association with the disliked is suffering, separation from the loved is suffering;}
  \chantlineinline
    {Yampicchaṃ na labhati, tampi dukkhaṃ.}
    {not getting what one wants is also suffering.}
\end{chantsubsection}
\chantgap
% ---------- 第二段 ----------
\begin{chantsubsection}
  \chantlineinline
    {Saṅkhittena pañcupādānakkhandhā dukkhā.}
    {In brief, the five aggregates subject to clinging are suffering.}
  \chantlineinline
    {Seyyathīdaṃ, rūpūpādānakkhandho vedanūpādānakkhandho
     saññūpādānakkhandho saṅkhārūpādānakkhandho
     viññāṇūpādānakkhandho,}
    {That is to say: the aggregate of clinging to form, to feeling, to perception,
     to formations, and to consciousness.}
  \chantlineinline
    {yesaṃ pariññāya dharamāno so bhagavā evaṃ bahulaṃ sāvake vineti.}
    {For the full understanding of these, while still alive, the Blessed One often
     trains his disciples in this way.}
  \chantlineinline
    {Evaṃbhāgā ca panassa bhagavato sāvakesu anusāsanī bahulā pavattati.}
    {And in this way his instruction among the disciples frequently proceeds.}
  \chantlineinline
    {Rūpaṃ aniccaṃ, vedanā aniccā, saññā aniccā, saṅkhārā aniccā,
     viññāṇaṃ aniccaṃ,}
    {Form is impermanent, feeling is impermanent, perception is impermanent,
     formations are impermanent, consciousness is impermanent;}
  \chantlineinline
    {rūpaṃ anattā, vedanā anattā, saññā anattā, saṅkhārā anattā,
     viññāṇaṃ anattā,}
    {form is not-self, feeling is not-self, perception is not-self,
     formations are not-self, consciousness is not-self;}
  \chantlineinline
    {sabbe saṅkhārā aniccā, sabbe dhammā anattāti.}
    {all conditioned things are impermanent, all dhammas are not-self.}
\end{chantsubsection}
\chantgap
% ---------- 第三段 ----------
\begin{chantsubsection}
  \chantlineinline
    {Te\footnote{For ladies, change \emph{Te} to \emph{Tā}.} mayaṃ, otiṇṇāmaha
     jātiyā jarāmaraṇena sokehi paridevehi dukkhehi domanassehi upāyāsehi.}
    {We, who are sunk in birth, aging and death, in sorrow, lamentation, pain,
     grief and despair,}
  \chantlineinline
    {Dukkhotiṇṇā dukkhaparetā, appevanāmimassa kevalassa
     dukkhakkhandhassa antakiriyā paññāyethāti.}
    {overwhelmed by suffering and oppressed by suffering, may the ending of this
     entire mass of suffering become apparent to us.}
  \chantlineinline
    {Ciraparinibbutampi taṃ bhagavantaṃ saraṇaṃ gatā,
     dhammañca saṅghañca,}
    {Though that Blessed One has long attained final Nibbāna, we have gone for
     refuge to that Blessed One, to the Dhamma and to the Saṅgha,}
  \chantlineinline
    {tassa bhagavato sāsanaṃ yathāsati yathābalaṃ manasikaroma
     anupaṭipajjāma.}
    {and we attend to that Blessed One’s teaching and practise it according to
     our mindfulness and strength.}
  \chantlineinline
    {Sā sā no paṭipatti imassa kevalassa dukkhakkhandassa
     antakiriyāya saṃvattatu.}
    {May this practice of ours lead to the ending of this entire mass of suffering.}
\end{chantsubsection}

%========================================
% Section：[9] Taṅkhaṇika-paccavekkhaṇa-pāṭho
%========================================
% title-en: Reflection at the Time of Use
\section*{m9 | Taṅkhaṇika-paccavekkhaṇa-pāṭho}

\begin{chantsubsection}
  \chantnote
    {\{Handa mayaṃ taṅkhaṇika-paccavekkhaṇa-pāṭhaṃ bhaṇāmase\}}
  \chantnotetrans
    {\{Now let us recite the reflection at the time of using the requisites.\}}
  \chantgap

  % Cīvara (robe)
  \chantlineinline
    {Paṭisaṅkhā yoniso cīvaraṃ paṭisevāmi,}
    {Wisely reflecting, I use the robe,}
  \chantlineinline
    {yāvadeva sītassa paṭighātāya,}
    {only for warding off cold,}
  \chantlineinline
    {uṇhassa paṭighātāya,}
    {for warding off heat,}
  \chantlineinline
      {ḍaṃsa-makasa-vātātapa-siriṃsapa-samphassānaṃ paṭighātāya,}
    {for warding off the touch of flies, mosquitoes, wind, sun, and creeping things,}
  \chantlineinline
    {yāvadeva hirikopina-paṭicchādanatthaṃ.}
    {only for the purpose of covering the body.}

\chantgap

  % Piṇḍapāta (almsfood)
  \chantlineinline
    {Paṭisaṅkhā yoniso piṇḍapātaṃ paṭisevāmi,}
    {Wisely reflecting, I use almsfood:}
  \chantlineinline
    {neva davāya,}
    {not for fun,}
  \chantlineinline
    {na madāya,}
    {not for intoxication,}
  \chantlineinline
    {na maṇḍanāya,}
    {not for fattening,}
  \chantlineinline
    {na vibhūsanāya,}
    {not for beautification,}
  \chantlineinline
    {yāvadeva imassa kāyassa ṭhitiyā, yāpanāya,}
    {but only for the maintenance and nourishment of this body,}
  \chantlineinline
    {vihiṃsūparatiyā, brahmacariyānuggahāya,}
    {for keeping it unharmed and for supporting the holy life, thinking:}
  \chantlineinline
    {iti purāṇañca vedanaṃ paṭihaṅkhāmi,}
    {'thus I shall allay old feeling,}
  \chantlineinline
    {navañca vedanaṃ na uppādessāmi,}
    {and not arouse new feeling;}
  \chantlineinline
    {yātrā ca me bhavissati,}
    {there will be for me support for the journey of life,}
  \chantlineinline
    {anavajjatā ca phāsuvihāro cā'ti.}
    {and blamelessness and comfortable abiding.}

\end{chantsubsection}
\chantgap
\begin{chantsubsection}

  % Senāsana (lodging)
  \chantlineinline
    {Paṭisaṅkhā yoniso senāsanaṃ paṭisevāmi,}
    {Wisely reflecting, I use a lodging,}
  \chantlineinline
    {yāvadeva sītassa paṭighātāya,}
    {only for warding off cold,}
  \chantlineinline
    {uṇhassa paṭighātāya,}
    {for warding off heat,}
  \chantlineinline
    {ḍaṃsa-makasa-vātātapa-siriṃsapa-samphassānaṃ paṭighātāya,}
    {for warding off the touch of flies, mosquitoes, wind, sun, and creeping things,}
  \chantlineinline
    {yāvadeva utuparissaya vinodanaṃ paṭisallānārāmatthaṃ.}
    {only for dispelling the hardships of climate and for the enjoyment of seclusion.}

\chantgap

  % Bhesajja (medicinal requisites)
  \chantlineinline
    {Paṭisaṅkhā yoniso gilāna-paccaya-bhesajja-parikkhāraṃ paṭisevāmi,}
    {Wisely reflecting, I use medicinal requisites for the sick,}
  \chantlineinline
    {yāvadeva uppannānaṃ veyyābādhikānaṃ vedanānaṃ paṭighātāya,}
    {only for warding off painful feelings that have arisen,}
  \chantlineinline
    {abyāpajjha-paramatāyā'ti.}
    {and for the attainment of freedom from affliction.}
\end{chantsubsection}

%========================================
% Section：[10] Dhātupaṭikūlapaccavekkhaṇapāṭha
%========================================
\section*{m10 | Dhātupaṭikūla-paccavekkhaṇa-pāṭho}

\begin{chantsubsection}
  \chantnote
    {\{Handa mayaṃ dhātupaṭikūlapaccavekkhaṇapāṭhaṃ bhaṇāma se\}}
  \chantnotetrans
    {\{Now let us recite the contemplation on the repulsiveness of the elements.\}}
\chantgap
  % Cīvara (robe)
  \chantlineinline
    {Yathāpaccayaṃ pavattamānaṃ dhātumattamevetaṃ,}
    {Arising according to conditions, this is merely a mass of elements,}
  \chantlineinline
    {yadidaṃ cīvaraṃ,}
    {namely this robe,}
  \chantlineinline
    {tadupabhuñjako ca puggalo dhātumattako nissatto nijjīvo suñño,}
    {and the person who makes use of it is also just elements, devoid of self, lifeless, empty,}
  \chantlineinline
    {sabbāni pana imāni cīvarāni ajigucchaniyāni}
    {yet all these robes themselves are not disgusting}
  \chantlineinline
    {imaṃ pūtikāyaṃ patvā ativiya jigucchaniyāni jāyanti.}
    {but, when applied to this putrid body, they become exceedingly disgusting.}
\end{chantsubsection}
\chantgap
\begin{chantsubsection}
  % Piṇḍapāta (almsfood)
  \chantlineinline
    {Yathāpaccayaṃ pavattamānaṃ dhātumattamevetaṃ,}
    {Arising according to conditions, this is merely a mass of elements,}
  \chantlineinline
    {yadidaṃ piṇḍapāto,}
    {namely this almsfood,}
  \chantlineinline
    {tadupabhuñjako ca puggalo dhātumattako nissatto nijjīvo suñño,}
    {and the person who eats it is also just elements, devoid of self, lifeless, empty,}
  \chantlineinline
    {sabbo panāyaṃ piṇḍapāto ajigucchaniyo}
    {yet this whole almsfood itself is not disgusting}
  \chantlineinline
    {imaṃ pūtikāyaṃ patvā ativiya jigucchaniyo jāyati.}
    {but, on reaching this putrid body, it becomes exceedingly disgusting.}
\end{chantsubsection}
\chantgap
\begin{chantsubsection}
  % Senāsana (lodging)
  \chantlineinline
    {Yathāpaccayaṃ pavattamānaṃ dhātumattamevetaṃ,}
    {Arising according to conditions, this is merely a mass of elements,}
  \chantlineinline
    {yadidaṃ senāsanaṃ,}
    {namely this lodging,}
  \chantlineinline
    {tadupabhuñjako ca puggalo dhātumattako nissatto nijjīvo suñño,}
    {and the person who uses it is also just elements, devoid of self, lifeless, empty,}
  \chantlineinline
    {sabbāni pana imāni senāsanāni ajigucchaniyāni}
    {yet all these lodgings themselves are not disgusting}
  \chantlineinline
    {imaṃ pūtikāyaṃ patvā ativiya jigucchaniyāni jāyanti.}
    {but, in connection with this putrid body, they become exceedingly disgusting.}
\end{chantsubsection}
\chantgap
\newpage
\begin{chantsubsection}
  % Bhesajja (medicinal requisites)
  \chantlineinline
    {Yathāpaccayaṃ pavattamānaṃ dhātumattamevetaṃ,}
    {Arising according to conditions, this is merely a mass of elements,}
  \chantlineinline
    {yadidaṃ gilānapaccayabhesajjaparikkhāro,}
    {namely medicines and requisites for the sick,}
  \chantlineinline
    {tadupabhuñjako ca puggalo dhātumattako nissatto nijjīvo suñño,}
    {and the person who uses them is also just elements, devoid of self, lifeless, empty,}
  \chantlineinline
    {sabbo panāyaṃ gilānapaccaya-bhesajjaparikkhāro ajigucchaniyo}
    {yet this whole support of medicines for the sick is itself not disgusting}
  \chantlineinline
    {imaṃ pūtikāyaṃ patvā ativiya jigucchaniyo jāyati.}
    {but, when it comes to this putrid body, it becomes exceedingly disgusting.}
\end{chantsubsection}

%========================================
% Section：[11] Pattidānagāthā
%========================================
\section*{m11 | Pattidāna-gāthā}

\begin{chantsubsection}
    \chantnote
    {\{Handa mayaṃ pattidāna gāthāyo bhaṇāma se\}}
    \chantnotetrans{\{Now let us recite the verses for sharing merit.\}}
\chantgap
  \chantline
    {Yā devatā santi vihāravāsinī, }
    {May the devas who dwell in this monastery,}
  \chantline
    {thūpe ghare bodhighare tahiṃ tahiṃ, }
    {in the stupa, the buildings, and the Bodhi-tree enclosure, here and there,}
  \chantline
    {tā dhammadānena bhavantu pūjitā, }
    {be honoured through this gift of Dhamma,}
  \chantline
    {sotthiṃ karontedha vihāramaṇḍale, }
    {and bring safety and well-being to this monastery compound.}

  \chantline
    {therā ca majjhā navakā ca bhikkhavo, }
    {May elder, intermediate, and newly ordained monks,}
  \chantline
    {sārāmikā dānapatī upāsakā, }
    {temple attendants, donors, and lay followers,}
  \chantline
    {gāmā ca desā nigamā ca issarā, }
    {villages, districts, towns, and their rulers,}
  \chantline
    {sappāṇabhūtā sukhitā bhavantu te, }
    {and all beings endowed with life, be happy.}

  \chantline
    {jalābujā yepi ca aṇḍasambhavā, }
    {Those born from the womb and those who come from an egg,}
  \chantline
    {saṃsedajātā athavopapātikā, }
    {those born from moisture, or who arise spontaneously,}
  \chantline
    {niyyānikaṃ dhammavaraṃ paṭicca te, }
    {having relied on the excellent Dhamma that leads out of suffering,}
  \chantline
    {sabbepi dukkhassa karontu saṅkhayaṃ, }
    {may they all bring about the ending of suffering.}

%% 下面开始 chantpair，节省空间
  \chantlinepair
    {thātu ciraṃ sataṃ dhammo, }
    {dhammaddharā ca puggalā, }
    {May the Dhamma of the good endure for a long time, together with the individuals who uphold the Dhamma.}

  \chantlinepair
    {saṅgho hotu samaggo va, }
    {atthāya ca hitāya ca, }
    {May the Sangha remain harmonious, for welfare and for benefit.}

  \chantlinepair
    {amhe rakkhatu saddhammo, }
    {sabbepi dhammacārino, }
    {May the true Dhamma protect us and all who practise in accordance with the Dhamma.}

  \chantlinepair
    {vuḍḍhiṃ sampāpuṇeyyāma, }
    {dhamme ariyappavedite. }
    {May we attain growth and fulfilment in the Dhamma taught and proclaimed by the Noble Ones.}

\end{chantsubsection}

% Vince proofread

%========================================
% Section [12] Ovāda-pāṭimokkha-gāthā
%========================================
\section*{m12 | Ovāda-pāṭimokkha-gāthā}

\begin{chantsubsection}
  \chantnote
    {\{Handa mayaṃ buddhassa ovādanusāsaniyo bhaṇāmase\}}
  \chantnotetrans
    {\{Now let us recite the Buddha's words of advice and instruction.\}}
\end{chantsubsection}
\chantgap
\begin{chantsubsection}
  % 第一段 Khantī...
  \chantnote
    {Khantī paramaṃ tapo titikkhā,}
  \chantindent
    {Patience is the supreme purifying practice,}
  \par\chantindent
    {(khuam-ot-thon khue khuam-ot-klan pen ta-ba yang-ying)}
  \chantgap

  \chantnote
    {Nibbānaṃ paramaṃ vadanti Buddhā,}
  \chantindent
    {Nibbāna is supreme, say the Buddhas.}
    \par
  \chantindent
    {(phra-phut-tha-jao tang-lai trat-wa phra-nip-phan pen-yiam)}
  \chantgap

  \chantnote
    {Na hi pabbajito parūpaghātī, samaṇo hoti paraṃ viheṭhayanto.}
  \chantindent
    {A mendicant does not harm others, a recluse oppresses no one.}
  \par\chantindent
    {(phu-lang-phlan phoo-uen mai-chue wa-pen ban-pha-chit, phoo-biat-bian phoo-uen mai-chue-wa pen-samana-loei)}
  \chantgap

  \chantnote
    {Etaṃ Buddhānaṃ sāsanaṃ.}
  \chantindent
    {This is the teaching of all Buddhas.}
    \par
  \chantindent
    {(nī-pen kham-sorn khong phra-phut-tha-jao tang-lai)}
  \chantgap

  % 第二段 Sabbapāpassa...
  \chantnote
    {Sabbapāpassa akaraṇaṃ,}
  \chantindent
    {Abstaining from all evil,}
  \par\chantindent
    {(kan-mai-tham-bap-tang-puang nueng)}
  \chantgap

  \chantnote
    {kusalassūpasampadā,}
  \chantindent
    {cultivating the wholesome,}
    \par
  \chantindent
    {(kan-bam-phen ku-son hai-thueng-phrom nueng)}
  \chantgap

  \chantnote
    {sacittapariyodapanaṃ,}
  \chantindent
    {cleansing of one's own mind,}
    \par
  \chantindent
    {(kan-klan-chit khong-ton hai-phong-phaew nueng)}
  \chantgap

  \chantnote
    {etaṃ Buddhānaṃ sāsanaṃ.}
  \chantindent
    {This is the teaching of all Buddhas.}
  \par\chantindent
    {(nī pen-kham-sorn khong phra-phut-tha-jao tang-lai)}
  \chantgap

  \clearpage
  
  % 第三段 Anūpavādo...
  \chantnote
    {anūpavādo,}
  \chantindent
    {not insulting,}
  \par\chantindent
    {(kan-mai-khao-pai wa-rai-kan nueng)}
  \chantgap

  \chantnote
    {anūpaghāto,}
  \chantindent
    {not harming,}
  \par\chantindent
    {(kan-kai-khao-pai lang-phlan-kan nueng)}
  \chantgap

  \chantnote
    {pāṭimokkhe ca saṃvaro,}
  \chantindent
    {carefulness in the Fundamental Precepts,}
  \par\chantindent
    {(khuam-sam-ruam nai phra-pa-ti-mok nueng)}
  \chantgap

  \chantnote
    {mattaññutā ca bhattasmiṃ,}
  \chantindent
    {moderation in food,}
  \par\chantindent
    {(khuam-pen phoo-ru-pra-man nai pho-cha-na-han nueng)}
  \chantgap

  \chantnote
    {pantañca sayanāsanaṃ,}
  \chantindent
    {dwelling in solitude,}
  \par\chantindent
    {(kan-non kan-nang an-sa-ngad nueng)}
  \chantgap

  \chantnote
    {adhicitte ca āyogo,}
  \chantindent
    {engaging in higher mental development,}
  \par\chantindent
    {(kan-pra-kop khuam-phian nai a-thi-jit nueng)}
  \chantgap

  \chantnote
    {etaṃ Buddhānaṃ sāsanaṃ.}
  \chantindent
    {This is the teaching of all Buddhas.}
  \par\chantindent
    {(nī pen-kham-sorn khong phra-phut-tha-jao tang-lai)}
\end{chantsubsection}

    
% \end{chantsubsection}

% 原文（需要校对字符）
% ```
% Ovada-patimokkha-gatha
% {Handa mayam buddhassa ovadanusasanlyo bhanama se.}

% KhantI paramam tapo titikkha
% Nibbanam paramam vadanti buddha
% Na hi pabbajito parupaghati
% Samano hoti param vihethayanto
% Etam buddhanasasanam.
% Sabbapapassa akaranam
% Kusalass'upasampadil
% Sacittapariyodapanam
% Etam buddhanasasanam.
% Anupavado anupaghato
% Patimokkhe ca samvaro
% Mattannuta ca bhattasmim
% Pan tan ca sayanasanam
% Adhicitte ca ayogo
% Etam buddhanasasanan'ti.
% ```

% 原文对应的英文和泰语翻译。
% ```
%  * Khantī paramaṃ tapo titikkhā,
%    Patience is the supreme purifying practice,
%    (khuam-ot-thon khue khuam-ot-klan pen ta-ba yang-ying)
%  * nibbānaṃ paramaṃ vadanti buddhā,
%    Nibbāna is supreme, say the Buddhas.
%    (phra-phut-tha-jao tang-lai trat-wa phra-nip-phan pen-yiam)
%  * na hi pabbājito parūpaghātī samaṇo hoti paraṃ viheṭhayanto,
%    A mendicant does not harm others, A recluse oppresses no one.
%    (phu-lang-phlan phoo-uen mai-chue-wa pen ban-pha-chit, phoo- biat-bian phoo-uen mai-chue-wa pen-samana-loei)
%  * etaṃ buddhāna sāsanaṃ.
%    This is the teaching of all Buddhas.
%    (ni pen-kham-sorn khong phra-phut-tha-jao tang-lai)
%  * sabbapāpassa akaraṇaṃ,
%    Abstaining from all evil,
%    (kan-mai-tham-bap-puang nueng)
%  * kusalassūpasampadā,
%    Doing all good thing,
%    (kan-bam-phen ku-son hai-thueng-phrom nueng)
%  * sacittapariyodapanaṃ,
%    Cleansing of one's mind,
%    (kan-klan-chit khong-ton hai-chue-phaew nueng)
%  * etaṃ buddhāna sāsanaṃ.
%    This is the teaching of all Buddhas.
%    (ni pen-kham-sorn khong phra-phut-tha-jao tang-lai)
%  * anūpavādo,
%    Not insulting,
%    (kan-mai-khao-pai wa-rai-kan nueng)
%  * anūpaghāto,
%    Not harming,
%    (kan-kai-khao-pai lang-phlan-kan nueng)
%  * pātimokkhe ca saṃvaro,
%    Carefulness of the Fundamental Precepts,
%    (khuam-sam-ruam nai phra-pa-ti-mok nueng)
%  * mattaññutā ca bhattasmiṃ,
%    Moderating in food,
%    (khuam-pen phoo-ru-pra-man nai pho-cha-na-han nueng)
%  * pantañca sayanāsanaṃ,
%    Dwelling in solitude,
%    (kan-non kan-nang an-sa-ngad nueng)
%  * adhicitte ca āyogo,
%    Engaging in higher mental development,
%    (kan-pra-kop khuam-phian nai a-thi-jit nueng)
%  * etaṃ buddhāna sāsanaṃ.
%    This is the teaching of all Buddhas.
%    (ni pen-kham-sorn khong phra-phut-tha-jao tang-lai)
% ```

% 我需要的格式举例。和之前的不同，三种语言直接放在一起，方便用多语种吟诵。
% ```
% \section*{[12] ...}

% \begin{chantsubsection}
%     \chantnote
%     {\{Handa mayaṃ pattidāna gāthāyo bhaṇāma se\}}
%     \chantindent{\{Now let us recite the verses for sharing merit.\}}
% \end{chantsubsection}

% \begin{chantsubsection}
%     \chantnote
%     {Khanti paramam tapo titikkha}
%     Patience is the supreme purifying
%     \par
%     (Khuam-ot-thon khue khuam-ot-klan pen ta-ba yang-ying)

%     \chantnote
%     ...
% \end{chantsubsection}
% ```
\newpage% Vince proofread

% title-en: Mettā Chanting (Spreading Loving Kindness)
\section*{m13 | Mettā Chanting (Spreading Loving Kindness)}\label{sec:m13}

\begin{chantsubsection}
  % Sabbe sattā – 总说一切众生
  \chantlineinline
    {Sabbe sattā}
    {All beings | All beings are our companions, born, growing old, falling ill, and dying together.}
  \chantindent
    {Sat-tang-lai, tee-ben-puen-tuk, gert-gae-jep-die, duay-gan mot-tang-sin.}
  \chantgap

  % Averā – 无瞋
  \chantlineinline
    {averā}
    {free from enmity | May we truly be happy and never use violence against one another.}
  \chantindent
    {Jong-ben-suk ben-suk-terd. Ya-dai-mee-wehn gae-gan-lae-gan-leuy.}
  \chantgap

  % Abyāpajjhā – 不互害
  \chantlineinline
    {abyāpajjhā}
    {free from ill will | May we truly be happy and never injure or oppress each other.}
  \chantindent
    {Jong-ben-suk-ben-suk-terd, ya-dai-biat-bian-seung-gan lae-gan-leuy.}
  \chantgap

  % Anīghā – 无苦恼
  \chantlineinline
    {anīghā}
    {free from affliction | May we truly be happy and have no suffering of body or mind at all.}
  \chantindent
    {Jong-ben-suk-ben-suk-terd, ya-dai-mee-kwam-tuk-gai-tuk-jai-leuy.}
  \chantgap

  % Sukhī attānaṃ pariharantu – 乐而自护
  \chantlineinline
    {sukhī attānaṃ pariharantu}
    {May they dwell in happiness | May you have happiness of body and mind, protecting yourselves from all suffering; may those who suffer find happiness, and may those who are happy grow even happier.}
  \chantindent
    {Jong-mee-kwam-suk-gai suk-jai, rak-sa-ton-hai-pon jak-tuk-pai, tang-sin-terd.}
  \chantindent
    {Tan-tang-lai, tee-tan dai-tuk, kor-hai-tan mee-kwam-suk, tan-tang-lai, tee-tan-dai-suk, kor-hai-suk-ying-ying-keurn-bai.}
  \chantgap

  % Sabbe sattā – 四生众生皆受功德
  \chantlineinline
    {sabbe sattā}
    {all beings | All beings who are egg-born, womb-born, moisture-born, or spontaneously born: may they receive this wholesome merit and blessing, each and every one.}
  \chantindent
    {Sat-tang-lai, tee-gerd-pen aṇ-ḍa-ja, tee-gerd-pen ja-lā-bu-ja, tee-gerd-pen saṃ-se-da-ja, tee-gerd-pen o-pa-pāti-ka, jong-ma-rap-gu-son-pohn boon, hai-tuan-tua tuk-tua-tuern.}
\end{chantsubsection}


% title-en: Salutation to the Triple Gem
% body-mode: bilingual
\section*{m14 | Ratanattayanāmakārapātha}

\begin{chantsubsection}
  % The Buddha
  \chantnote
    {Arahaṃ sammāsambuddho bhagavā, buddhaṃ bhagavantaṃ abhivādemi.}
  \chantnotetrans
    {The Blessed One is an arahant, perfectly awakened; to the Buddha, the Blessed One, I pay homage.}
  \chantgap
  
  \chantnote
    {--- bow down and chant softly ---}
  \chantnote
    {“buddho me nātho.}
  \chantnote
    {The Buddha is my refuge.”}
  \chantindent
    {(phra-phut-tha-jao pen thi-phueng khong-rao)}
  \chantgap
  \chantgap

  % The Dhamma
  \chantnote
    {Svākkhāto bhagavatā dhammo, dhammaṃ namassāmi.}
  \chantnotetrans
    {The Dhamma is well expounded by the Blessed One; to the Dhamma I pay homage.}
  \chantnote
    {--- bow down and chant softly ---}
  \chantnote
    {“dhammo me nātho.}
  \chantnote
    {The Dhamma is my refuge.”}
  \chantindent
    {(phra-dham pen thi-phueng khong-rao)}
  \chantgap
  \chantgap
  
  % The Saṅgha
  \chantnote
    {Supaṭipanno bhagavato sāvaka-saṅgho, saṅghaṃ namāmi.}
  \chantnotetrans
    {The Sangha of the Blessed One's disciples has practiced well; to the Sangha I pay homage.}
  \chantnote
    {--- bow down and chant softly ---}
  \chantnote
    {“saṅgho me nātho.}
  \chantnote
    {The Saṅgha is my refuge.”}
  \chantindent
    {(phra-song pen thi-phueng khong-rao)}
\end{chantsubsection}

% 结尾说明（可选）
\chantgap
\chantgap
{\englishstyle
This saying shows respect for the Triple Gem, firms faith in it, and keeps a close relationship with it in mind. It serves as a preliminary step that will make the attainment of Dhammakaya (the Body of Truth) more accessible.}

\chantgap
\chantgap
{\englishstyle\textit{End of the morning chanting}
}


\newpage

\chanttitle{Evening Chanting}

% 01-ratanattaya.tex
% Section [1] Ratanattaya-vandanā
% title-en: Homage to the Triple Gem

\section*{e1 | Ratanattaya-vandanā}

\begin{chantsubsection}
  \chantline
    {Yo so bhagavā arahaṃ sammāsambuddho}
    {For translation, see page~\pageref{sec:m01} (Morning Chanting)}

  \chantonly
    {Svākkhāto yena bhagavatā dhammo}

  \chantonly
    {Supaṭipanno yassa bhagavato sāvakasaṅgho}

  \chantonly
    {Tammayaṃ bhagavantaṃ sadhammaṃ sasaṅghaṃ}

  \chantonly
    {Imehi sakkārehi yathārahaṃ āropitehi abhipūjayāma}

  \chantonly
    {Sādhu no bhante bhagavā suciraparinibbutopi}

  \chantonly
    {Pacchimājanatānukampamānasā}

  \chantonly
    {Ime sakkāre duggatapaṇṇākārabhūte paṭiggaṇhātu}

  \chantonly
    {Amhākaṃ dīgharattaṃ hitāya sukhāya}
\end{chantsubsection}


%========================================
% Section：[2] Ratanattayanamakārapāṭha
%========================================
% title-en: Salutation to the Triple Gem
\section*{e2 | Ratanattayanamakārapāṭha}\label{sec:m02}

\begin{chantsubsection}
  \chantline
    {Arahaṃ sammāsambuddho bhagavā}
    {For translation, see page~\pageref{sec:m02} (Morning Chanting)}

  \chantonly
    {Buddhaṃ bhagavantaṃ abhivādemi}
  \chantnote{---bow---}

  \chantonly
    {Svākkhāto bhagavatā dhammo dhammaṃ namassāmi}

  \chantnote{---bow---}

  \chantonly
    {Supaṭipanno bhagavato sāvakasaṅgho saṅghaṃ namāmi}
    
  \chantnote{---bow---}

\end{chantsubsection}

\newpage\input{e0304-Pubbabhaganamakarapatha}
%========================================
% [5] Buddhābhigīti
%========================================
% title-en: Chant in Praise of the Buddha
\section*{e5 | Buddhābhigīti}

\begin{chantsubsection}
  \chantnote
    {\{Handa mayaṃ buddhābhigītiṃ karomase\}}
  \chantnotetrans
    {\{Now let us chant in praise of the Buddha.\}}
\end{chantsubsection}

% ---- Opening verse ----
\begin{chantsubsection}
  \chantline
    {Buddhavārahantavaratādiguṇābhiyutto}
    {Endowed with the noble qualities of the Buddha, the Arahant, and the Excellent Ones,}
  \chantline
    {Suddhābhiñāṇakaruṇāhi samāgatatto}
    {his being complete with pure insight and compassion,}
  \chantline
    {Bodhesi yo sujanataṃ kamalaṃ va sūro}
    {he awakened the good people, like the sun [awakens] the lotus,}
  \chantline
    {Vandāmahaṃ taṃ araṇaṃ sirasā jinendaṃ}
    {I bow with my head to that peaceful one, the Lord of Conquerors.}
\end{chantsubsection}

% ---- Main praise (two pādas per line) ----
\begin{chantsubsection}
  \chantlinepair
    {Buddho yo sabbapāṇīnaṃ}
    {Saraṇaṃ khemamuttamaṃ}
    {The Buddha is the secure, supreme refuge for all living beings.}

  \chantlinepair
    {Paṭhamānussatiṭṭhānaṃ}
    {Vandāmi taṃ sirenāhaṃ}
    {As the first object of recollection, I bow to him with my head.}

  % footnote 1: use \footnotemark inside minipage text, then \footnotetext right after
  \chantlinepair
    {Buddhassāhaṃ dāso\footnotemark{} va}
    {Buddho me sāmikissaro}
    {I am the Buddha's servant; the Buddha is my lord and master.}
  \footnotetext{For ladies: change \textit{dāso} to \textit{dāsī}.}

  \chantlinepair
    {Buddho dukkhassa ghātā ca}
    {Vidhātā ca hitassa me}
    {The Buddha destroys suffering and brings about my welfare.}

  \chantlinepair
    {Buddhassāhaṃ niyyādemi}
    {Sarīrañjīvitañcidaṃ}
    {To the Buddha I dedicate this body and this very life.}

  % footnote 2
  \chantlinepair
    {Vandantohaṃ\footnotemark{} carissāmi}
    {Buddhasseva subodhitaṃ}
    {Paying homage, I will live (practice) in accordance with the Buddha’s clearly realized Dhamma.}
  \footnotetext{For ladies: change \textit{Vandantohaṃ} to \textit{Vandantīhaṃ}.}

  \chantlinepair
    {Natthi me saraṇaṃ aññaṃ}
    {Buddho me saraṇaṃ varaṃ}
    {I have no other refuge; the Buddha is my supreme refuge.}

  \chantlinepair
    {Etena saccavajjena}
    {Vaḍḍheyyaṃ satthu sāsane}
    {By this truthful utterance, may I grow in the Teacher’s dispensation.}

  % footnote 3
  \chantlinepair
    {Buddhaṃ me vandamānena\footnotemark{}}
    {Yaṃ puññaṃ pasutaṃ idha}
    {By my paying homage to the Buddha, whatever merit is produced here,}
  \footnotetext{For ladies: change \textit{vandamānena} to \textit{vandamānāya}.}

  \chantlinepair
    {Sabbepi antarāyā me}
    {Māhesuṃ tassa tejasā}
    {may all obstacles not afflict me through the power of his radiance.}
\end{chantsubsection}

\newpage

\begin{chantsubsection}
  \chantnote{\textit{--- bow, chanting softly ---}}
\end{chantsubsection}

% ---- Confession verse ----
\begin{chantsubsection}
  \chantline
    {Kāyena vācāya va cetasā vā}
    {By body, speech, or mind,}
  \chantline
    {Buddhe kukammaṃ pakataṃ mayā yaṃ}
    {whatever wrong I have done toward the Buddha,}
  \chantline
    {Buddho paṭiggaṇhatu accayantaṃ}
    {may the Buddha accept my confession,}
  \chantline
    {Kālantare saṃvarituṃ va buddhe}
    {so that in the future I may restrain myself with regard to the Buddha.}
\end{chantsubsection}

%========================================
% 6 | Dhammānussati (short)
%========================================
\section*{e6 | Dhammānussati}

\begin{chantsubsection}
  \chantnote{\{Handa mayaṃ dhammānussatinayaṃ karomase\}}
  \chantnotetrans{\{Now let us chant in recollection of the Dhamma.\}}
\end{chantsubsection}

\begin{chantsubsection}
  \chantlineinline
    {Svākkhāto bhagavatā dhammo,}
    {The Dhamma is well-expounded by the Blessed One,}
  \chantlineinline
    {sandiṭṭhiko akāliko ehipassiko,}
    {visible here and now, timeless, inviting one to come and see,}
  \chantlineinline
    {opanayiko paccattaṃ veditabbo viññūhī’ti.}
    {leading inward, to be personally known by the wise.}
\end{chantsubsection}

%========================================
% Section：[7] Dhammābhigīti
%========================================
% title-en: Chant in Praise of the Dhamma
\section*{e7 | Dhammābhigīti}

\begin{chantsubsection}
  \chantline
    {Handa mayaṃ dhammābhigītiṃ karoma se}
    {Now let us recite a chant in praise of the Dhamma.}

  \chantline
    {Svākkhātatādiguṇayogavasena seyyo}
    {The Dhamma is supreme, endowed with such qualities as being well expounded,}

  \chantline
    {Yo maggapākapariyattivimokkhabhedo}
    {distinguished as Path, Fruit, Study, and Liberation.}

  \chantline
    {Dhammo kulokapatanā tadadhāridhārī}
    {The Dhamma prevents the fall into the bad world and supports those who uphold it.}

  \chantline
    {Vandāmahaṃ tamaharaṃ varadhammametaṃ}
    {I revere that excellent Dhamma which dispels darkness.}

  \chantlinepair
    {Dhammo yo sabbapāṇīnaṃ}
    {Saraṇaṃ khemamuttamaṃ}
    {The Dhamma is for all beings the supreme, secure refuge.}

  \chantlinepair
    {Dutiyānussatiṭṭhānaṃ}
    {Vandāmi taṃ sirenahaṃ}
    {It is the second foundation of recollection; to it I pay homage with my head.}

  % footnote 1
  \chantlinepair
    {Dhammassāhasmi dāso\footnotemark{} va}
    {Dhammo me sāmikissaro}
    {I am a servant of the Dhamma; the Dhamma is my lord and master.}
  \footnotetext{For ladies: change \textit{dāso} to \textit{dāsī}.}

  \chantlinepair
    {Dhammo dukkhassa ghātā ca}
    {Vidhātā ca hitassa me}
    {The Dhamma is the destroyer of suffering and the establisher of my welfare.}

  \chantlinepair
    {Dhammassāhaṃ niyyādemi}
    {Sarīrañjīvitañcidaṃ}
    {To the Dhamma I entrust this body and this life.}

  % footnote 2
  \chantlinepair
    {Vandantohaṃ\footnotemark{} carissāmi}
    {Dhammasseva sudhammataṃ}
    {Paying homage, I will live devoted only to the true nature of the Dhamma.}
  \footnotetext{For ladies: change \textit{Vandantohaṃ} to \textit{Vandantīhaṃ}.}

  \chantlinepair
    {Natthi me saraṇaṃ aññaṃ}
    {Dhammo me saraṇaṃ varaṃ}
    {I have no other refuge; the Dhamma is my excellent refuge.}

  \chantlinepair
    {Etena saccavajjena}
    {Vaḍḍheyyaṃ satthu sāsane}
    {By this truthful utterance, may I grow in the Teacher's dispensation.}

  % footnote 3
  \chantlinepair
    {Dhammaṃ me vandamānena\footnotemark{}}
    {Yaṃ puññaṃ pasutaṃ idha}
    {Whatever merit I have gathered here by revering the Dhamma,}
  \footnotetext{For ladies: change \textit{vandamānena} to \textit{vandamānāya}.}

  \chantlinepair
    {Sabbepi antarāyā me}
    {Māhesuṃ tassa tejasā}
    {may all my obstacles be dispelled through its power.}
\end{chantsubsection}

%========================================
% 8 | Saṅghānussati
%========================================
% title-en: Recollection of the Sangha
\section*{e8 | Saṅghānussati}

\begin{chantsubsection}
  \chantnote
    {\{Handa mayaṃ saṅghānussatinayaṃ karomase\}}
  \chantnotetrans
    {\{Now let us chant in recollection of the Saṅgha.\}}
\end{chantsubsection}

\begin{chantsubsection}
  \chantline
    {Supaṭipanno bhagavato sāvakasaṅgho}
    {The Saṅgha of the Blessed One’s disciples has practiced well.}

  \chantline
    {Ujupaṭipanno bhagavato sāvakasaṅgho}
    {The Saṅgha of the Blessed One’s disciples has practiced uprightly.}

  \chantline
    {Ñāyapaṭipanno bhagavato sāvakasaṅgho}
    {The Saṅgha of the Blessed One’s disciples has practiced rightly.}

  \chantline
    {Sāmīcipaṭipanno bhagavato sāvakasaṅgho}
    {The Saṅgha of the Blessed One’s disciples has practiced properly.}

  \chantline
    {Yadidaṃ cattāri purisayugāni aṭṭha purisapuggalā}
    {That is to say, the four pairs of persons, the eight kinds of individuals.}

  \chantline
    {Esa bhagavato sāvakasaṅgho}
    {This is the Saṅgha of the Blessed One’s disciples.}

  \chantline
    {Āhuneyyo pāhuneyyo dakkhiṇeyyo añjalīkaraṇīyo}
    {Worthy of gifts, worthy of hospitality, worthy of offerings, worthy of reverential salutation.}

  \chantline
    {Anuttaraṃ puññakkhettaṃ lokassāti}
    {The unsurpassed field of merit for the world.}
\end{chantsubsection}

%========================================
% 9 | Saṅghābhigīti
%========================================
% title-en: Chant in Praise of the Sangha
\section*{e9 | Saṅghābhigīti}

\begin{chantsubsection}
  \chantnote
    {\{Handa mayaṃ saṅghābhigītiṃ karomase\}}
  \chantnotetrans
    {\{Now let us recite a chant in praise of the Saṅgha.\}}
\end{chantsubsection}

\chantgap

% ---- Opening praise ----
\begin{chantsubsection}
  \chantline
    {Saddhammajo supaṭipattiguṇādiyutto}
    {Born of the True Dhamma, endowed with virtues such as good practice,}

  \chantline
    {Yoṭṭhabbidho ariyapuggalasaṅghaseṭṭho}
    {the noble Saṅgha, supreme among the eight kinds of noble persons,}

  \chantline
    {Sīlādidhammapavarāsayakāyacitto}
    {with body and mind established in the excellent abode of virtue and other qualities,}

  \chantline
    {Vandāmahaṃ tamariyānagaṇaṃ susuddhaṃ}
    {I revere that pure and noble assembly.}
\end{chantsubsection}

\newpage

% ---- Main praise (two pādas per line) ----
\begin{chantsubsection}
  \chantlinepair
    {Saṅgho yo sabbapāṇīnaṃ}
    {Saraṇaṃ khemamuttamaṃ}
    {The Saṅgha is for all beings the supreme, secure refuge.}

  \chantlinepair
    {Tatiyānussatiṭṭhānaṃ}
    {Vandāmi taṃ sirenahaṃ}
    {As the third foundation of recollection, I bow to it with my head.}

  % footnote 1
  \chantlinepair
    {Saṅghassāhasmi dāso\footnotemark{} va}
    {Saṅgho me sāmikissaro}
    {I am a servant of the Saṅgha; the Saṅgha is my lord and master.}
  \footnotetext{For ladies: change \textit{dāso} to \textit{dāsī}.}

  \chantlinepair
    {Saṅgho dukkhassa ghātā ca}
    {Vidhātā ca hitassa me}
    {The Saṅgha destroys suffering and brings about my welfare.}

  \chantlinepair
    {Saṅghassāhaṃ niyyādemi}
    {Sarīrañjīvitañcidaṃ}
    {To the Saṅgha I dedicate this body and this very life.}

  % footnote 2
  \chantlinepair
    {Vandantohaṃ\footnotemark{} carissāmi}
    {Saṅghassopaṭipannataṃ}
    {Paying homage, I will live in accordance with the Saṅgha’s good practice.}
  \footnotetext{For ladies: change \textit{Vandantohaṃ} to \textit{Vandantīhaṃ}.}

  \chantlinepair
    {Natthi me saraṇaṃ aññaṃ}
    {Saṅgho me saraṇaṃ varaṃ}
    {I have no other refuge; the Saṅgha is my excellent refuge.}

  \chantlinepair
    {Etena saccavajjena}
    {Vaḍḍheyyaṃ satthu sāsane}
    {By this truthful utterance, may I grow in the Teacher’s dispensation.}

  % footnote 3 (note: first pāda stands alone in the source)
  \chantlinepair
    {Saṅghaṃ me vandamānena\footnotemark{}}
    {Yaṃ puññaṃ pasutaṃ idha}
    {By my paying homage to the Saṅgha, whatever merit is produced here,}
  \footnotetext{For ladies: change \textit{vandamānena} to \textit{vandamānāya}.}

  \chantlinepair
    {Sabbepi antarāyā me}
    {Māhesuṃ tassa tejasā}
    {may all my obstacles be dispelled through its power.}
\end{chantsubsection}

% ---- Gesture note ----
\chantgap
\begin{chantsubsection}
  \chantnote
    {\textit{--- bow, chanting softly ---}}

\chantgap

% ---- Confession verse ----
  \chantline
    {Kāyena vācāya va cetasā vā}
    {By body, speech, or mind,}

  \chantline
    {Saṅghe kukammaṃ pakataṃ mayā yaṃ}
    {whatever wrong I have done toward the Saṅgha,}

  \chantline
    {Saṅgho paṭiggaṇhatu accayantaṃ}
    {may the Saṅgha accept my confession,}

  \chantline
    {Kālantare saṃvarituṃ va saṅghe}
    {so that in the future I may restrain myself with regard to the Saṅgha.}
\end{chantsubsection}

%========================================
% 11 | Uddissanagāthā
%========================================
\section*{e11 | Uddissanagāthā}

\begin{chantsubsection}
  \chantnote
    {\{Handa mayaṃ uddissanagāthāyo bhaṇāmase\}}
  \chantnotetrans
    {\{Now let us recite the verses of dedication.\}}
\end{chantsubsection}

\begin{chantsubsection}
  \chantlinepair
    {Iminā puññakammena}
    {Upajjhāyā guṇuttarā}
    {By this meritorious deed, may my preceptors of superior virtue [benefit].}

  \chantlinepair
    {Ācariyūpakārā ca}
    {Mātā pitā ca ñātakā piyā mamaṃ}
    {May my teachers and benefactors, my mother and father, and my dear relatives [benefit].}

  \chantlinepair
    {Suriyo candimā rājā}
    {Guṇavantā narāpi ca}
    {The sun, the moon, kings, and virtuous people as well,}

  \chantlinepair
    {Brahmamārā ca indā ca}
    {Lokapālā ca devatā}
    {the Brahmās, Māras, Indas, guardian deities, and gods,}

  \chantlinepair
    {Yamo mittā manussā ca}
    {Majjhattā verikāpi ca}
    {Yama, friends, human beings, the impartial, and even enemies,}

  \chantlinepair
    {Sabbe sattā sukhī hontu}
    {Puññāni pakatāni me}
    {may all beings be happy through the merits I have made.}

  \chantlinepair
    {Sukhaṃ ca tividhaṃ dentu}
    {Khippaṃ pāpetha vomataṃ}
    {May they bestow the threefold happiness and quickly reach Nibbāna.}

  \chantlinepair
    {Iminā puññakammena}
    {Iminā uddisena ca}
    {By this meritorious deed and by this dedication,}

  \chantlinepair
    {Khippāhaṃ sulabhe ceva}
    {Taṇhupādānachedanaṃ}
    {may I swiftly and easily attain the cutting off of craving and clinging.}

  \chantlinepair
    {Ye santāne hinā dhammā}
    {Yāva nibbānato mamaṃ}
    {Whatever inferior states exist in my mental continuum until Nibbāna,}

  \chantlinepair
    {Nassantu sabbadā yeva}
    {Yattha jāto bhave bhave}
    {may they always disappear, wherever I am reborn, life after life.}

  \chantlinepair
    {Ujucittaṃ satipaññā}
    {Sallekho vīriyamhinā}
    {May I have a straight mind, mindfulness and wisdom, and purification through effacement and energy.}

  \chantlinepair
    {Mārā labhantu nokāsaṃ}
    {Kātuñca vīriyesu me}
    {May Māra find no opportunity to obstruct my efforts.}

  \chantlinepair
    {Buddho dīpavaro nātho}
    {Dhammo nātho varuttamo}
    {The Buddha, the excellent lamp, is my refuge; the Dhamma is my supreme refuge.}

  \chantlinepair
    {Nātho paccekabuddho ca}
    {Saṅgho nāthottaro mamaṃ}
    {The Paccekabuddha is my refuge, and the Saṅgha is my highest refuge.}

  \chantlinepair
    {Tesottamānubhāvena}
    {Mārokāsaṃ labhantu mā}
    {By the power of these supreme ones, may Māra gain no opportunity [over me].}
\end{chantsubsection}

%========================================
% 12 | Abhiṇha-paccavekkhaṇā-pāṭho
%========================================
\section*{e12 | Abhiṇha-paccavekkhaṇā-pāṭho}\label{sec:e12}

\begin{chantsubsection}
  \chantnote
    {\{Handa mayaṃ abhiṇhapaccavekkhaṇapāṭhaṃ bhaṇāmase\}}
  \chantnotetrans
    {\{Now let us recite the Daily Reflections.\}}

  \chantgap

  \chantnote
    {Jarādhammomhi jarāṃ anatīto\footnote{For gentlemen: \textit{anatīto}; for ladies: \textit{anatītā}.}}
  \chantnote
    {I am of the nature to age; I have not gone beyond aging.}
  \chantindent
    {(rao mi khuam-kae pen tham-ma-da yang-mai luang-phon khuam-kae pai-dai)}

  \chantgap

  \chantnote
    {Byādhidhammomhi byādhiṃ anatīto\footnote{For gentlemen: \textit{anatīto}; for ladies: \textit{anatītā}.}}
  \chantnote
    {I am of the nature to be ill; I have not gone beyond illness.}
  \chantindent
    {(rao mi khuam-jep pen tham-ma-da yang-mai luang-phon khuam-jep pai-dai)}

  \chantgap

  \chantnote
    {Maraṇadhammomhi maraṇaṃ anatīto\footnote{For gentlemen: \textit{anatīto}; for ladies: \textit{anatītā}.}}
  \chantnote
    {I am of the nature to die; I have not gone beyond death.}
  \chantindent
    {(rao mi khuam-tai pen tham-ma-da yang-mai luang-phon khuam-tai pai-dai)}

  \chantgap

  \chantnote
    {Sabbehi me piyehi manāpehi nānābhāvo vinābhāvo}
  \chantnote
    {All that is mine, dear and delightful, will change and be separated from me.}
  \chantindent
    {(rao ja-tong phlat-phrak chak khong-rak khong-chop-jai duai-kan-mod-tang-sin)}

  \chantgap

  \chantnote
    {Kammassakomhi kammadāyādo}
  \chantnote
    {I am the owner of my kamma; the heir to my kamma.}
  \chantindent
    {(rao mi-kam pen khong-tua pen tha-yat haeng-kam)}

  \chantgap

  \chantnote
    {Kammayoni kammabandhu}
  \chantnote
    {Born of my kamma; related to my kamma.}
  \chantindent
    {(mi kam-pen-kam-noet mi-kam-pen-phao-phan)}

  \chantgap
  \clearpage

  \chantnote
    {Kammapaṭisaraṇo\footnote{For gentlemen: \textit{saraṇo}; for ladies: \textit{saraṇā}.}}
  \chantnote
    {I rely on my kamma.\footnote{To rely on kamma: beyond all external supports,
one must return to and bear the results of one's own intentional actions.}}
  \chantindent
    {(mi kam pen-thi-phueng ar-sai)}

  \chantgap

  \chantnote
    {Yaṃ kammaṃ karissāmi}
  \chantnote
    {Whatever kamma I shall do.}
  \chantindent
    {(rao tham-kam-dai-wai)}

  \chantgap

  \chantnote
    {Kalyāṇaṃ vā pāpakaṃ vā}
  \chantnote
    {Whether good or evil,}
  \chantindent
    {(dee-rue-chua ko-tam)}

  \chantgap

  \chantnote
    {Tassa dāyādo\footnote{For gentlemen: \textit{dāyādo}; for ladies: \textit{dāyādā}.} bhavissāmi}
  \chantnote
    {Of that I shall be the heir.}
  \chantindent
    {(rao ja-tong pen-phoo-rap phon-khong-kam-nan)}
    
\end{chantsubsection}

\end{document}
```

参考资源二
``` morning chanting section 10
%========================================
% Section：[10] Dhātupaṭikūlapaccavekkhaṇapāṭha
%========================================
\section*{10 | Dhātupaṭikūla-paccavekkhaṇa-pāṭho}

\begin{chantsubsection}
  \chantnote
    {\{Handa mayaṃ dhātupaṭikūlapaccavekkhaṇapāṭhaṃ bhaṇāma se\}}
  \chantnotetrans
    {\{Now let us recite the contemplation on the repulsiveness of the elements.\}}
\chantgap
  % Cīvara (robe)
  \chantlineinline
    {Yathāpaccayaṃ pavattamānaṃ dhātumattamevetaṃ,}
    {Arising according to conditions, this is merely a mass of elements,}
  \chantlineinline
    {yadidaṃ cīvaraṃ,}
    {namely this robe,}
  \chantlineinline
    {tadupabhuñjako ca puggalo dhātumattako nissatto nijjīvo suñño,}
    {and the person who makes use of it is also just elements, devoid of self, lifeless, empty,}
  \chantlineinline
    {sabbāni pana imāni cīvarāni ajigucchaniyāni}
    {yet all these robes themselves are not disgusting}
  \chantlineinline
    {imaṃ pūtikāyaṃ patvā ativiya jigucchaniyāni jāyanti.}
    {but, when applied to this putrid body, they become exceedingly disgusting.}
\end{chantsubsection}
\chantgap
\begin{chantsubsection}
  % Piṇḍapāta (almsfood)
  \chantlineinline
    {Yathāpaccayaṃ pavattamānaṃ dhātumattamevetaṃ,}
    {Arising according to conditions, this is merely a mass of elements,}
  \chantlineinline
    {yadidaṃ piṇḍapāto,}
    {namely this almsfood,}
  \chantlineinline
    {tadupabhuñjako ca puggalo dhātumattako nissatto nijjīvo suñño,}
    {and the person who eats it is also just elements, devoid of self, lifeless, empty,}
  \chantlineinline
    {sabbo panāyaṃ piṇḍapāto ajigucchaniyo}
    {yet this whole almsfood itself is not disgusting}
  \chantlineinline
    {imaṃ pūtikayaṃ patvā ativiya jigucchaniyo jāyati.}
    {but, on reaching this putrid body, it becomes exceedingly disgusting.}
\end{chantsubsection}
\chantgap
\begin{chantsubsection}
  % Senāsana (lodging)
  \chantlineinline
    {Yathāpaccayaṃ pavattamānaṃ dhātumattamevetaṃ,}
    {Arising according to conditions, this is merely a mass of elements,}
  \chantlineinline
    {yadidaṃ senāsanaṃ,}
    {namely this lodging,}
  \chantlineinline
    {tadupabhuñjako ca puggalo dhātumattako nissatto nijjīvo suñño,}
    {and the person who uses it is also just elements, devoid of self, lifeless, empty,}
  \chantlineinline
    {sabbāni pana imāni senāsanāni ajigucchaniyāni}
    {yet all these lodgings themselves are not disgusting}
  \chantlineinline
    {imaṃ pūtikāyaṃ patvā ativiya jigucchaniyāni jāyanti.}
    {but, in connection with this putrid body, they become exceedingly disgusting.}
\end{chantsubsection}
\chantgap
\begin{chantsubsection}
  % Bhesajja (medicinal requisites)
  \chantlineinline
    {Yathāpaccayaṃ pavattamānaṃ dhātumattamevetaṃ,}
    {Arising according to conditions, this is merely a mass of elements,}
  \chantlineinline
    {yadidaṃ gilānapaccayabhesajjaparikkhāro,}
    {namely medicines and requisites for the sick,}
  \chantlineinline
    {tadupabhuñjako ca puggalo dhātumattako nissatto nijjīvo suñño,}
    {and the person who uses them is also just elements, devoid of self, lifeless, empty,}
  \chantlineinline
    {sabbo panāyaṃ gilānapaccaya-bhesajjaparikkhāro ajigucchaniyo}
    {yet this whole support of medicines for the sick is itself not disgusting}
  \chantlineinline
    {imaṃ pūtikāyaṃ patvā ativiya jigucchaniyo jāyati.}
    {but, when it comes to this putrid body, it becomes exceedingly disgusting.}
\end{chantsubsection}
```

根据参考，处理下面内容
要求：检查原文拼写、翻译、编排

```
10 ｜ Atīta-paccavekkhaṇa-vidhī
{Handa mayaṃ atītapaccavekkhaṇapāṭhaṃ bhaṇāmase.}

Ajja mayā apaccavekkhitvā yaṃ cīvaraṃ paribhuttaṃ, taṃ yāvadeva sītassa paṭighātāya, uṇhassa paṭighātāya, ḍaṃsamakasavātātapasiriṃsapasamphassānaṃ paṭighātāya, yāvadeva hirikopinapaṭicchādanatthaṃ.

Ajja mayā apaccavekkhitvā yo piṇḍapāto paribhutto, so neva davāya na madāya na maṇḍanāya na vibhūsanāya, yāvadeva imassa kāyassa ṭhitiyā yāpanāya vihiṃsuparatiyā brahmacariyānuggahāya, iti purāṇañca vedanaṃ paṭihaṅkhāmi navañca vedanaṃ na uppādessāmi, yātrā ca me bhavissati anavajjatā ca phāsuvihāro cāti.

Ajja mayā apaccavekkhitvā yaṃ senāsanaṃ paribhuttaṃ, taṃ yāvadeva sītassa paṭighātāya, uṇhassa paṭighātāya, ḍaṃsamakasavātātapasiriṃsapasamphassānaṃ paṭighātāya, yāvadeva utuparissayavinodanaṃ paṭisallānārāmatthaṃ.

Ajja mayā apaccavekkhitvā yo gilānapaccayabhesajjaparikkhāro paribhutto, so yāvadeva uppannānaṃ veyyābādhikānaṃ vedanānaṃ paṭighātāya, abyāpajjhaparamatāyāti.
```
